\section{引言}
\label{sec:introduction}

在 UTXO 账本中，一笔支付花费的是先前某笔交易创建的输出。如果收款方必须等到该交易完成公共执行，支付链的每一跳都会增加一次确认延迟。设想 Alice 向 Bob 支付 100 CAL，Bob 再把收到的输出转给 Carol。要在公共执行之前转出这个输出，Bob 需要可验证的证据，证明其价值有资金支撑。如果 Bob 的支付先执行，账本就消费了一个 Alice 的交易尚未创建的输出，缺失的 100 CAL 需要一个明确的资金来源。无论 Alice 的交易先于 Bob 的到达、后于它到达，还是始终不到达，账务都必须保持正确。

基于证书的支付系统已经把收款方一侧的可用性与公共处理分离开来。FastPay 和 Zef 支持低延迟的认证支付，Sui Lutris 则针对不同类型的状态访问，把广播路径与共识路径结合起来~\cite{fastpay,zef,lutris}。支付通道协议依靠通道关系及其流动性实现快速转账~\cite{blitz,donner}。Next 处理的是原生账本上已认证支付的资金支撑问题，特别是所消费输出的源交易尚未执行的情形。源交易尚未执行的已认证输出，应当能够支撑一笔已公共执行的后继交易，并且应当由一个已提供备付资金的明确责任方对缺失的源交易负责。用户花费指定了消费路由的自有输出，注册组织为未解决的源交易提供支撑，交易对手无需预先注资的通道路径。

这里存在四个相互耦合的问题。第一，证书在公共登记之前就已流通，因此从公开登记的缺口开始记账会漏掉已经存在的承诺，而签名额度必须覆盖每一笔已认证支付的未清偿 CAL 负债，包括从未公开提交的支付。第二，源交易可能在委员会已经为其缺失的价值赔付之后才到达。迟到的源交易不能把同一笔价值既付给垫付它的担保人，又付给已经花掉它的收款方，否则延迟本身就会转移资金。第三，客户端可能持有源交易的材料却从不提交。签发组织的成员应能在赔付到期前补交已有的源交易，而赔付无需等待历史适配服务。第四，源交易一旦被赔付，读取已存储历史的人应当能够看出已消费的输入由哪份备付金提供资金。记录这一事实必须保持所有者授权、输出标识符以及后继交易所引用的承诺不变。仅有稳定的交易哈希，既无法界定也无法授权这样的修改。

Next 固定支付权利，并把未解决的源交易划归其直接签发组织负责。输出证书（TXCer）把一笔支付的输出绑定到一个明确的签发组织。收款方在验证并原子地记录输出及其 TXCer 之后，通过该输出指定的消费组织请求后继交易。该请求携带其直接输入的证书。所有者授权、花费状态、费用和资源限额则另行检查。

如果某个已认证的输入在公共创建状态中仍然缺失，委员会就原子地消费它，记录针对签发组织的直接义务，并创建子输出。在义务处于 Open 状态期间，子输出即为最终输出。承载该子交易的区块同时承载创建了被消费输出的那笔支付的证书。批准过该支付的签发组织成员，根据自己保存的批准记录重建该支付并补交，无需新的签名。如果源交易执行，义务即告履行。如果在期限之后赔付决定执行时源交易仍未出现，该决定就扣减签发组织的备付金，由它提供缺失的支撑，收款方不会得到第二次入账。该决定无需门限适配。之后到达的源交易在同一转换中偿还赔付它的签发组织，并且不为收款方创建第二个输出。当所有源交易最终都执行时，延迟只改变备付金的临时占用，本金的最终分配与无延迟时相同。

赔付决定本身就是一条关闭义务的公共记录。随后，历史表示命令在既有承诺之下，把已消费输入的资金引用替换为备付金引用。已存储的区块指明为该输入提供资金的备付金扣减，并且仍可对照其原始区块身份验证，因此读者能够在区块本身中核对指定历史输入的资金来源。钱包、同步和重放继续读取原始命令。我们把这一设计归纳为三项职责各异的贡献。

\noindent\textbf{C1：面向一轮续花的分层架构。}
我们把成员认证、收款方交付、证据复制和公共执行分离开来。成员在签名之前，原子地锁定输入并预留额度。一个并行认证轮产生一份 TXCer。INSTALL 在组织内复制完整材料，公共提交独立推进，二者均在后台进行。注册组织在组织内部和组织之间使用同一个证据接口，支付本金与担保人备付金保持分离。

\noindent\textbf{C2：有资金支撑且延迟中性的直接责任记账。}
我们定义原子的缺失输入转换及其关闭路径：由源交易履行、由公共赔付决定赔付、由迟到的源交易偿还。对于固定的输入与证书形态，请求的祖先证据不随转账深度增长。在固定成员的拜占庭模型下，我们证明唯一消费与资产守恒，并证明签名额度约束了证书（包括未公开提交的证书）的未清偿 CAL 负债。对于有限的、互不冲突且源交易全部执行的支付集合，到达与赔付的任意交错，最终得到的 CAL 持有和备付金余额都与无延迟执行相同。

\noindent\textbf{C3：证据驱动的源交易解决与保持身份的历史表示。}
我们让原签署者依据后继交易公开的证书补交源交易，并将赔付决定与门限适配分离。受限的历史表示命令没有经济效果，保持所有者授权、输出引用及完整区块身份（含分片集合头）不变。同一区块的命令共享受影响分片最终表示的适配，安装则直接推进到已提交的修订版本。对于固定的已决定槽位集合，最终区块体不依赖于命令如何分组。

我们在一台 M4 Max 主机上评估两个版本：续花与记账基线版本，以及增加了补交和独立赔付决定的 decision-first 版本。在基线版本上，100 跳链完成最后一次认证接收的三轮用时中位数为 158.846\,ms，而在各跳之间等待公共确认则需要 65.204\,s。三次各含 30,000 笔支付的运行（含后台排空）完成速率为 1,750--1,887 笔/秒，一个九次运行的矩阵关闭 21,600 笔支付，并偿还全部 216 次赔付。在 decision-first 版本上，三次运行的功能测试显示，被扣留的源交易无需赔付即可执行，而收款方的接收约需 14\,ms。在所有适配服务暂停的情况下，赔付与偿还仍然完成；六轮普通路径比较中，运行级接收 P50 和 P95 的范围相互重叠。第~\ref{sec:model} 节和第~\ref{sec:protocol} 节定义该设计，第~\ref{sec:security} 节和第~\ref{sec:evaluation} 节展开其安全性论证与评估。
